%----------------------------------------------------------------------------------------
%	CURRICULUM VITAE EXECUTIVO — HUGO SOARES (pdfLaTeX, coluna única + cabeçalho)
%----------------------------------------------------------------------------------------
\documentclass[10pt,a4paper]{extarticle}

% --- Geometria da página -------------------------------------------------------
\usepackage[a4paper, left=0.55in, right=0.5in, top=0in, bottom=0.45in]{geometry}

% --- Codificação e tipos de letra ----------------------------------------------
\usepackage[utf8]{inputenc}
\usepackage[T1]{fontenc}
\usepackage{lmodern}
\usepackage{microtype}
\renewcommand{\familydefault}{\sfdefault}
\linespread{1.08}

% --- Paleta de cores ------------------------------------------------------------
\usepackage{xcolor}
\definecolor{primary}{HTML}{1A365D}    % Azul-marinho
\definecolor{secondary}{HTML}{2B6CB0}  % Azul-ardósia
\definecolor{ice}{HTML}{C7D9EF}        % Azul claro (sobre navy)
\definecolor{accent}{HTML}{E3EAF3}     % Fundo azul-claro
\definecolor{textdark}{HTML}{2D3748}   % Carvão
\definecolor{graylight}{HTML}{718096}  % Cinzento frio
\color{textdark}

% --- Pacotes --------------------------------------------------------------------
\usepackage{tikz}
\usepackage{enumitem}
\usepackage{array}
\usepackage{fontawesome5}
\usepackage{hyperref}

\hypersetup{
    colorlinks=true,
    linkcolor=secondary,
    urlcolor=secondary,
    pdftitle={Curriculum Vitae — Hugo Soares},
    pdfauthor={Hugo Soares},
}

% --- Controlo de hifenização -----------------------------------------------------
\tolerance=1
\emergencystretch=\maxdimen
\hyphenpenalty=10000
\hbadness=10000

% --- Estilo das listas -----------------------------------------------------------
\setlist[itemize]{
    leftmargin=1.15em,
    topsep=1.5pt,
    itemsep=1.6pt,
    parsep=0pt,
    partopsep=0pt,
    label=\textcolor{secondary}{\textbullet},
    before=\color{textdark},
}

% ==============================================================================
%	MACROS
% ==============================================================================
% Cabeçalho de secção (maiúsculas, navy, espaçado, sublinhado)
\newcommand{\cvh}[1]{%
    \par\vspace{5pt}%
    \noindent{\color{primary}\large\bfseries\textls[130]{\MakeUppercase{#1}}}\par
    \vspace{0pt}%
    \noindent{\color{secondary}\rule{\linewidth}{1pt}}\par
    \vspace{2pt}%
}

% Separador de contacto
\newcommand{\sep}{{\color{ice}\normalsize\raisebox{0.1ex}{\textbullet}}\hspace{2pt}}

% Linha de cabeçalho de projeto: título (+github) à esquerda, datas à direita
\newcommand{\projhead}[3]{%
    \noindent
    \begin{tabular}{@{}>{\bfseries\color{primary}}l@{\extracolsep{\fill}}>{\scriptsize\color{graylight}\raggedleft\arraybackslash}r@{}}
        #1 & #2%
    \end{tabular}\par
    \vspace{0.5pt}%
    {\scriptsize\color{secondary}\textbf{Stack:} #3}\par
    \vspace{2pt}%
}

% ==============================================================================
%	INÍCIO DO DOCUMENTO
% ==============================================================================
\begin{document}
\pagestyle{empty}

% ==============================================================================
%	FAIXA DE CABEÇALHO FULL-BLEED (navy em toda a largura, com foto)
% ==============================================================================
\begingroup
\setlength{\fboxsep}{0pt}
\hypersetup{urlcolor=ice}
\noindent\makebox[0pt][l]{%
\hspace*{-0.55in}%
\colorbox{primary}{%
\begin{minipage}[c]{\paperwidth}%
    \vspace{0.22cm}%
    \hspace*{0.55in}%
    \begin{minipage}[c]{\dimexpr\paperwidth-0.55in-0.5in-2.9cm-0.5cm\relax}%
        {\color{white}\fontsize{20}{22}\selectfont\bfseries Hugo Soares}%
        {\color{ice}\normalsize\bfseries\ |\ Engenheiro de Software Full-Stack}\par
        \vspace{4pt}
        {\small\color{white}
            \faMapMarker\ Braga, Portugal \sep
            \faEnvelope\ \href{mailto:hugoferreirasoares@gmail.com}{hugoferreirasoares@gmail.com}\par
            \vspace{2pt}
            \faLinkedin\ \href{https://www.linkedin.com/in/yhugosoares/}{linkedin.com/in/yhugosoares} \sep
            \faGithub\ \href{https://github.com/yhugosoares}{github.com/yhugosoares} \sep
            \faPhone\ +351 934 120 467}
    \end{minipage}%
    \hfill
    \begin{minipage}[c]{2.9cm}%
        \raggedleft
        \begin{tikzpicture}
            \begin{scope}
                \clip (0,0) circle (1.35cm);
                \node[anchor=center] at (0,0) {\includegraphics[width=2.7cm,height=2.7cm]{square_photo_compressed.jpg}};
            \end{scope}
            \draw[white,line width=1.4pt] (0,0) circle (1.35cm);
        \end{tikzpicture}
    \end{minipage}%
    \hspace*{0.5in}%
    \vspace{0.22cm}%
\end{minipage}}
}% close makebox
\endgroup

\vspace{6pt}

% ==============================================================================
%	RESUMO / PERFIL
% ==============================================================================
\cvh{Perfil}
{\small\color{textdark}
Licenciado em Engenharia Informática pela Universidade do Minho que concebe e entrega aplicações completas de ponta a ponta. Linguagens principais Java e TypeScript, com experiência prática na integração de serviços de IA, Docker e CI/CD em projetos entregues de forma independente em \textbf{Desktop}, \textbf{Web} e \textbf{Mobile}.}

% ==============================================================================
%	EDUCAÇÃO
% ==============================================================================
\cvh{Educação}
\noindent
\begin{tabular}{@{}>{\bfseries\color{primary}}l@{\extracolsep{\fill}}>{\scriptsize\color{graylight}\raggedleft\arraybackslash}r@{}}
    Licenciatura em Engenharia Informática | Universidade do Minho & 2023 -- 2026
\end{tabular}

\begin{itemize}
    \item Média: 14/20 valores.
    \item UCs relevantes: Programação Orientada a Objetos (Java), Análise e Teste de Software (JUnit, EvoSuite), Segurança de Sistemas Informáticos, Redes de Computadores, Processamento de Linguagens.
\end{itemize}

% ==============================================================================
%	COMPETÊNCIAS TÉCNICAS
% ==============================================================================
\cvh{Competências Técnicas}
\noindent
\begin{tabular}{@{}>{\color{secondary}\small\bfseries}p{3.3cm}@{}>{\small\color{textdark}\raggedright\arraybackslash}p{\dimexpr\linewidth-3.3cm\relax}@{}}
    Linguagens:       & Java, TypeScript, JavaScript, Python, Rust, C\\[2.5pt]
    Frameworks:       & React, Node.js, Express, Tauri\\[2.5pt]
    Bases de Dados:   & PostgreSQL, Supabase, Git, Linux/Bash\\[2.5pt]
    Cloud \& DevOps:  & Docker, Docker Compose, GitHub Actions (CI/CD)\\[2.5pt]
    IA \& LLMs:       & DeepSeek API, Ollama, LLMs Locais, Prompt Engineering\\[2.5pt]
    Testes:           & JUnit, EvoSuite, Hypothesis, Property-Based Testing\\[2.5pt]
    Idiomas:          & Português (Nativo), Inglês (Fluente)
\end{tabular}

% ==============================================================================
%	PROJETOS
% ==============================================================================
\cvh{Projetos}

\projhead{Reddit Story Pipeline \href{https://github.com/yHugoSoares/reddit-tiktok-pipeline}{\ \footnotesize\faGithub}}{2025 -- Ago 2026}{Python, DeepSeek API, Docker, Automação}
\begin{itemize}
    \item Concebeu e entregou um pipeline automatizado end-to-end que transforma histórias do Reddit em vídeos verticais narrados e legendados, integrando a API DeepSeek (LLM) para reescrita, voz TTS e legendas.
    \item Estruturou uma arquitetura containerizada com Docker Compose, gestão persistente de sessões e publicação multiplataforma para TikTok, Instagram e YouTube.
    \item Construiu fiabilidade desde o início: persistência de estado, deduplicação de publicações, modo dry-run para rollout faseado e agendamento via cron/Ofelia para até 3 publicações por dia.
\end{itemize}

\projhead{Zenyth IDE \href{https://github.com/yHugoSoares/Zenyth-IDE}{\ \footnotesize\faGithub}}{Nov 2025 -- Presente}{React, TypeScript, Tauri, Rust, Ollama}
\begin{itemize}
    \item Desenvolveu um IDE desktop nativo que transmite I/O de terminal bidirecionalmente sem perda de dados, permitindo execução de código em tempo real a partir de uma interface React envolta em Tauri.
    \item Projetou um proxy em Rust para inferência local de LLM (Ollama) com latência abaixo de 100ms no primeiro token, gerindo backpressure, cancelamento e recuperação graciosa de erros.
    \item Arquitetou a ponte IPC entre React e Rust que porta um editor web para uma aplicação desktop totalmente nativa.
\end{itemize}

\projhead{VencordMobile \href{https://github.com/yHugoSoares/VencordMobile}{\ \footnotesize\faGithub}}{2025 -- Presente}{Dart (Flutter), JavaScript, WebView, CI/CD}
\begin{itemize}
    \item Entregou um cliente Android que carrega a app web móvel do Discord e injeta o bundle de plugins do Vencord (155+ plugins), habilitando funcionalidades indisponíveis em qualquer outro cliente móvel.
    \item Desenvolveu uma ponte Flutter-JavaScript via postMessage para comunicação bidirecional nativo-web com sincronização de estado.
    \item Automatizou a entrega de ponta a ponta: bundling de 155+ plugins num único IIFE com esbuild, e GitHub Actions a produzir APKs assinados em cada release.
\end{itemize}

\end{document}
